\section{Conjecture de Goldbach}\label{sec:dm2-goldbach}

\begin{td-exo}[Preuve informelle d'une somme de points]\,\\
  On se donne des \(n\) points sur une droite. Décrire un algorithme
  qui détermine le plus petit sous-ensemble de points dont la somme
  des longueurs d'un segment est \(k\).
\end{td-exo}

\begin{td-sol}[]
  On construit un arbre de recherche, puis on procède par énumération
  du plus court chemin vers la racine.
\end{td-sol}

\section{Thème d'apprentissage : Programmation dynamique}\label{sec:dm2-dp}

\begin{td-exo}[Problème de bin packing]\,\\
  Soient \(B\) bins et \(I\) items d'intérêt \(w_1, \dots, w_n\).
  Expliquer comment construire un programme dynamique en temps polynomial
  en \(|I|\) et \(B\).
\end{td-exo}

\begin{td-sol}[]
  On définit \(OPT(i, b)\) comme le plus grand gain d'un subset de
  l'ensemble \(I\) à \(i\) éléments dans le bin \(b\).
  \(OPT(i, b) = \max\{ OPT(i-1, b), \; w_i + OPT(i-1, b - w_i) \}\).
\end{td-sol}
